%\documentclass{beamer}
%
%\usepackage[utf8]{inputenc}
%\usepackage[french]{babel}
%
%\begin{document}
%
%\title{Deductive Verification of Stateful Second-Order Programs}
%\author{Léon Gondelman}
%\date{M2 internship \\ \today}
%
%\begin{frame}
%  \maketitle
%\end{frame}


\documentclass[xcolor=dvipsnames]{beamer}
\usepackage{etex}
\usetheme{Pittsburgh}
\usepackage{lmodern} \linespread{1.1}
\usepackage[latin1]{inputenc}
\usepackage[sc]{mathpazo} \linespread{1.1}
\usepackage{changepage}
 \usepackage{float}
\usepackage{setspace}
\usepackage{pstricks,pstricks-add,pst-math,pst-xkey}
\usepackage{multicol}
\usepackage{bcprules, proof}

%\usepackage[T1]{fontenc}
\usepackage{graphicx}
\usepackage{hyperref}
\usepackage{listings}
\usepackage{mathdots}
\usepackage[usenames,dvipsnames]{}
\usepackage{amssymb,amsfonts,amsthm,amsmath}
\RequirePackage{listings}
\RequirePackage{amssymb}
\usepackage{tikz}

\lstset{
	xleftmargin=\parindent,
  basicstyle={\ttfamily},
%  framesep=2pt,
%  frame=single,
  keywordstyle={\color{blue}},
  stringstyle=\itshape,
  commentstyle=\itshape,
  columns=[l]fullflexible,
  showstringspaces=false,
  mathescape=true
}

%\lstset{basicstyle={\ttfamily}}
\lstdefinelanguage{whycode}{
morekeywords=[1]{},%
string=[b]",%
showstringspaces=false,%
stringstyle=\itshape,%
commentstyle= \color{gray}\small\texttt, %
columns=[l]fullflexible,%
sensitive=true,%
morecomment=[s]{(*}{*)},%
escapeinside={*?}{?*},%
keepspaces=true}
\lstset{language=whycode,
emph={namespace,result,predicate,module,fun,while,done,function,inductive,type,use,clone,%
import,export,theory,end,in,match,with,do,let,ghost,ensures,requires,ensures,invariant, rec,variant,then,if,else,%
goal,axiom,lemma,forall, induction},
emphstyle=\color{blue}}

\lstset{language=whycode}
\lstnewenvironment{whycode}{\lstset{language=whycode}}{}

\label{whygray}
\lstdefinelanguage{whygray}{
morekeywords=[1]{},%
string=[b]",%
showstringspaces=false,%
stringstyle=\itshape,%
commentstyle= \color{gray}\small\texttt, %
columns=[l]fullflexible,%
sensitive=true,%
morecomment=[s]{(*}{*)},%
escapeinside={*?}{?*},%
keepspaces=true}

\lstnewenvironment{whygray}{
\lstset{
language=whygray, 
stringstyle=\color{gray},
emphstyle=\color{gray}}\gray}{\gray}





\let\why\lstinline
%\newcommand{\blue[1]}{\color{blue}}
\let\emph\alert
\beamertemplatenavigationsymbolsempty
\newcommand{\HRule}{\rule{\linewidth}{0.1mm}}

\newenvironment{aquote}[1]
  {\savebox\mybox{#1}\begin{quote}}
  {\signed{\usebox\mybox}\end{quote}}

\addtobeamertemplate{footline}{\hfill\insertframenumber/\inserttotalframenumber}


\newcommand{\rouge}[1]{\textcolor{red}{#1}}
\newcommand{\bwedge}{\boldsymbol{~\wedge~}}
\newcommand{\bvee}{\boldsymbol{~\vee~}}
\newcommand{\brarr}{\boldsymbol{~\Rightarrow~}}
\setbeamertemplate{frametitle}{\vspace*{1em}\hfill\insertframetitle} 
\addtocounter{framenumber}{-1}

\newcommand{\gb}{\beta}
\newcommand{\gba}[1]{\beta_{#1}}
\newcommand{\gbb}{\bot_\beta}
\newcommand{\gbt}{\top_\beta}

\newcommand{\gbr}{\textcolor{red}{\gb}}
\newcommand{\gbra}[1]{\textcolor{red}{\gba{#1}}}
\newcommand{\gbbr}{\textcolor{red}{\gbb}}
\newcommand{\gbtr}{\textcolor{red}{\gbt}}
\newcommand{\gbran}[1]{\textcolor{red}{\neg\gba{#1}}}
\newcommand{\bth}{\bot_\theta}
\newcommand{\brh}{\bot_\rho}  
\newcommand{\tth}{\top_\theta}
\newcommand{\trh}{\top_\rho}   

\newcommand{\ghosttyping}[6]{\vdash_{gh}~#1~:~#2,~(#3,~#4,~\rouge{#5})~\textcolor{blue}{\boldsymbol{#6}}}

\newcommand{\crond}[2]{
\tikz \node[draw,circle,scale=0.21,fill=#2]{\black \Huge  \textbf{#1}};}
\newcommand{\rond}[1]{
\tikz \node[draw,circle,scale=0.21,fill=red]{\black \Huge  \textbf{#1}};}
\begin{document}
\section*{Title}
\begin{frame}
\begin{center}
\vspace{-1.cm}
%\HRule

\textsc{\blue{\Large Deductive Verification of \\[0.2cm] 
Stateful Second-Order Programs }}\\
%\HRule
 \textsc{ \scriptsize master's degree second year internship}\\[2cm]
%{\small  L\'eon \textsc{Gondelman}}  \\[0.5cm]
%{\small Advisor: Jean-Christophe \textsc{Filli\a^tre}}
\bigskip
\begin{minipage}{0.3\textwidth}
\begin{flushleft} 
\emph{Author:}\\
L\'eon \textsc{Gondelman}
\end{flushleft}
\end{minipage}
\begin{minipage}{0.5\textwidth}
\begin{flushright} 
\emph{Supervisor:} \\
Jean-Christophe \textsc{Filli\^atre}
\end{flushright}
\end{minipage}

\vfill

% Bottom of the page
{\small \today}



%\includegraphics[scale=0.15]{image_title3}
\end{center}
\end{frame}
%%%%%%%%%%%%%%%%%%%%%%%%%%%%%%%%%%%%%%%%%%%%%%%%%%%%%%%%%%%%%%%%%%%%%%%%%%%%%%%%



% 1. contexte, problème, exemple introductif (array_iter)
% sum_loop et sum_iter en parallèle


\begin{frame}{formal verification informally}
\begin{quotation}
 ''The most important property of a program is whether it accomplishes the intention of its user.''\\
\hfill{\blue{C.A.R. Hoare}}
\end{quotation}
\bigskip
\end{frame}

\begin{frame}
\frametitle{the art of program verification}
\includegraphics{definition_en.mps}
\bigskip\pause
  \begin{enumerate}
  \item {\red specify} an (existing) program \pause

  % \begin{center}
  %   \includegraphics{approach3.mps}
  % \end{center}
  \item {\red generate} verification conditions \pause
  \item {\red discharge} them, using
    \begin{itemize}
    \item ATP when possible
    \item proof assistants otherwise
    \end{itemize}
  \end{enumerate}
\end{frame}

\begin{frame}
\frametitle{state of art}
{\red last 40 years}
	\begin{itemize}
		\item[] formal verification of imperative programs 
		\item[] first-order verification plateforms (Why3, Boogie, ...)
		\item[] emergence of efficient ATPs (Alt-Ergo, Z3, CVC4, ...)
	\end{itemize}\vspace*{1em} \pause
{\red more recently}
	\begin{itemize}
		\item[] higher-order functional programs (Coq, Pangoline,...)
		\item[] higher-order {\red stateful} programs (Who, F-star, ...)
	\end{itemize}
\pause
\bigskip
{\red challenge}: proof automation
\end{frame}

\section*{Introduction}
\subsection*{Computing the sum of array elements}
\begin{frame}[fragile]\frametitle{Computing the sum of array elements}
\begin{adjustwidth}{-0.3cm}{1.0cm}
\begin{small}
\begin{minipage}[t]{0.4\textwidth}
\begin{whycode}  
let sum_loop (a: int array)
= let s = ref 0 in
  let n = Array.length a in
  let i = ref 0 in
  while !i < n do
    s := !s + a.(!i);
    i := !i + 1
  done;
  !s
\end{whycode}
\end{minipage}\hfill
\begin{minipage}[t]{0.4\textwidth}
\begin{whycode} 
let sum_iter (a: int array)
= let s = ref 0 in
  Array.iter 
    (fun x -> s := !s + x) a;
  !s
\end{whycode}	
\end{minipage}
\end{small}
\end{adjustwidth}
\end{frame}

\subsection*{Verifying Sum\_loop}
\subsubsection*{Code}
\newcommand{\OliveGreen}[1]{\textcolor{OliveGreen}{#1}}
\newcommand{\Rouge}[1]{\textcolor{red}{#1}}
\newcommand{\Noir}[1]{\textcolor{black}{#1}}
\begin{frame}[fragile]
\frametitle<1>
{sum example with {loop}}
\frametitle<2>
{\quad program specification}
\frametitle<3>
{\quad loop invariant}
\frametitle<4>
{precondition}
\frametitle<5>
{array bound check}
\frametitle<6>
{preservation}
\frametitle<7>
{postcondition}
\vspace*{-3em}
\begin{displaymath} 	
\begin{array}{@{\hspace*{0em}}l@{\hspace*{1.5em}}ll@{\hspace*{1em}}l}

\only<-3>{\vspace*{2em}}
\only<4>{
 \rond{vc1}	 
 &\textcolor{red}{\emptyset}
 &\textcolor{red}{\vdash}
 &\textcolor{red}{(0 \leq 0 \bwedge 0 \leq n) \bwedge 0 = \sum_{~0\leq j < 0} a[j] \phantom{ phantom of the opera is here right now}}}

\only<5>{
	\rond{vc2}
	&\textcolor{red}{\mathcal{I} \bwedge i < n}
	&\textcolor{red}{ \vdash} 
	&\textcolor{red}{ 0 \leq i < n} \phantom{phantom of the opera is here right now}}


\only<6>{
	\rond{vc3}
	&\textcolor{red}{\mathcal{I} \bwedge i < n}
	&\textcolor{red}{\vdash} 
	&\textcolor{red}{s + a[i] = \sum_{~0 \leq j < i+1} a[j] \bwedge (0 \leq (i+1) \leq n)   \phantom{phantom of the opera is here right now}}}


\only<7>{
	\rond{vc4}
	&\textcolor{red}{ \mathcal{I} \bwedge i \geq n} 
	&\textcolor{red}{ \vdash}
	&\textcolor{red}{ s = \sum_{~0 \leq j < n} a[j] \phantom{phantom of the opera is here right now}}}

\end{array}
\end{displaymath}
\vspace*{-2em}
\begin{normalsize}
	\begin{whycode}  
let sum_loop (a: int array) *?\only<2->{\\ 
	{\blue ensures} $\{\texttt{ {\blue result} = }{\red~\sum_{~0 \leq j < \text{a.length}} a[j]}~\}$}?* 
= let s = ref 0 in
  let n = Array.length a in
  let i = ref 0 in
  *?\only<4>{\hspace*{-2em}\crond{vc1}{red}~}\only<5->{\hspace*{-2em}\crond{vc1}{gray}~}?*while !i < n do *?
  
  \only<3->{\\    
  ${ \blue \hspace*{1em} \texttt{invariant}}  
   ~\{{\red~\mathcal{I}: \texttt{!s = } \sum_{0\leq j < !i} a[j] \bwedge 0 \leq !i \leq \texttt{n}}~\}$} 
  ?*
	 	*?
	 	\only<5>{\hspace*{-2em}\crond{vc2}{red}~}
	 	\only<4,6->{\hspace*{-2em}\crond{vc2}{gray}~}?*  s := !s + a.(!i);  
    i := !i + 1
 *?\only<6>{\hspace*{-1.5em}\crond{vc3}{red}\hspace*{-0.0em}}
 	\only<4-5,7->{\hspace*{-1.5em}\crond{vc3}{gray}\hspace*{-0.0em}}?* done;
 *?\only<4-6>{\hspace*{-1.5em}\crond{vc4}{gray}}\only<7->{\hspace*{-1.5em}\crond{vc4}{red}}?* !s 
\end{whycode}\vspace*{-0.5em}
\end{normalsize}
\end{frame}

\subsubsection*{Conclusion}
\begin{frame}
\frametitle{automated proof}
\begin{itemize}
\item All verification conditions {\tiny \rond{vc1} - \rond{vc4}} are first-order formulas.
\item Automated theorem provers easily prove them. 
%\item<4-> {\color{red} Can we generate first-order VCs from that specification ?}
\end{itemize}
\bigskip\pause
\begin{center}
 {\color{red} Can we prove \texttt{sum\_iter} in a similar way ?}
\end{center}
\end{frame}

\subsection*{Verifying Sum\_iter}
\subsubsection*{Code}


\begin{frame}[fragile]
\frametitle<1>{sum example with {iterator}}
\frametitle<2>{program specification}
\frametitle<3>{invariant formula}
\begin{whycode}
 let sum_iter (a: int array) *?\only<2->{\\
	{\blue \texttt{ ensures }}\{\texttt{ {\blue result} =} 
	${\red ~\sum_{~0 \leq j < \text{a.length}} a[j]~}$\}}?* 
 = let s = ref 0 in
   Array.iter*?
   \only<3->{\\\hspace*{3.5em}
   {\red(fun i -> !s = $\sum_{~0\leq j < i} a[j]$)}\\}?* 
       (fun x ->  s := !s + x) a; 
     !s 
\end{whycode}
\end{frame}

\begin{frame}[fragile]
\frametitle<1>{iterator specification}
\frametitle<2>{additional parameter}
\frametitle<3>{loop initial and final states}
\frametitle<4>{loop preservation}
\frametitle<5-6>{pros and cons}
\frametitle<7>{an existing solution}
\frametitle<8>{pragmatic approach}
\begin{whycode}  
let iter *?\only<2->{{\red \texttt{(inv:int -> prop) }}}?*(f:int -> unit) (b: int array) *?\only<3->{\\ 
\texttt{\blue requires }\{{\red \texttt{ inv 0 }}\} \\
\texttt{\blue ensures~  }\{{\red \texttt{ inv b.length }}\}}?* 
= let rec aux (i: int) *?\only<4->{\\
\texttt{~~\blue requires }\{{\red \texttt{ inv i} $\bwedge (0 \leq i \leq b.length)$}        \} \\
\texttt{~~\blue ensures~   }\{{\red \texttt{ inv b.length }}\}}?*  
  = if i < Array.length b 
    then (f b[i] ; aux (i + 1)) 
  in aux 0
\end{whycode} 
\bigskip
\only<5>{
$$ \fbox{\texttt{{\blue pros}: specification above is {\blue modular}}} $$}
\only<6>{
$$ \fbox{\texttt{{\red cons}: specification above is {\red partial}}} $$}
\only<7>{
 $$\texttt{\OliveGreen{complete the specification}}$$\vspace*{-1.5em} 
$$ \textcolor{purple}{\texttt{may be difficult to prove automatically}}$$
}
\only<8>{
$$ \texttt{\OliveGreen{\underline{inline} the definition of iterator}} $$
\vspace*{-1.5em} 
$$ \hspace*{0.5em}\texttt{\OliveGreen{\underline{instantiate} functional arguments}}$$}
\end{frame}


%
%\subsubsection*{Problem}
%\begin{frame}
%\frametitle<1-3>{\quad  Yes, but ...}
%\frametitle<4->{\quad Pragmatic approach:}
%\begin{itemize}
%\item[-]<1-> function \texttt{Array\_iter} is second-order code;
%\item[-]<2-> invariant \texttt{inv} is a second-order formula;
%\item[-]<3-> {\color{red} we cannot generate first-order VCs.}\\[1cm]
%%\item[-]<4-> { Yes, by second-order inlining !}\\[1em]
%\only<4->{ A simple solution:}~\\[0.5em]
%\begin{itemize}
%\item[\color{blue} 1)]<4-> {\color{blue} inline the definition of array\_iter};\\[1em]
%\item[\color{blue} 2)]<4-> {\color{blue} instantiate functional parameters \texttt{f} and \texttt{inv} }.
%\end{itemize}
%\end{itemize}
%\end{frame}
%

\subsubsection*{Solution:Second-Order Inlining}
\begin{frame}[fragile]
%\frametitle{\OliveGreen{inline definition of iterator}}
\vspace*{0em}
\begin{footnotesize}
\begin{whycode} 
module Array 
... 
   let iter (inv:int -> prop) (f:int -> unit) (b: int array)
   requires { inv 0 } 
   ensures  { inv b.length } 
   = let rec aux (i: int)
     requires { inv i$\bwedge 0 \leq \texttt{i} \leq \texttt{b.length}$ }
     ensures  { inv b.length }
     = if i < Array.length b
       then (f b[i] ; aux (i + 1)) 
     in aux 0
...
end

let sum_iter (a: int array) 
ensures { $\texttt{result = }\sum_{~0 \leq j < \text{a.length}} a[j]$ }
= let s = ref 0 in
  Array.iter 
    (fun i -> !s = $\sum_{~0\leq j < i} a[j]$) (fun x -> s := !s + x) a; 
  !s
\end{whycode} 
\end{footnotesize}
\end{frame}
\addtocounter{framenumber}{-1}




\begin{frame}[fragile]
\frametitle{\OliveGreen{inline iter}}
\vspace*{0em}
\begin{footnotesize}
\begin{whygray}
let sum_iter (a: int array) 
ensures { $\texttt{result = }\sum_{~0 \leq j < \text{a.length}} a[j]$ }
= let s = ref 0 in 
\end{whygray}\vspace*{-1.15em} 
\begin{whycode}
  let iter (inv:int -> prop) (f:int -> unit) (b: int array)
  requires { inv 0 } 
  ensures  { inv b.length } 
  = let rec aux (i: int)
    requires { inv i$\bwedge 0 \leq \texttt{i} \leq \texttt{b.length}$ }
    ensures  { inv b.length }
    = if i < Array.length b
      then (f b[i] ; aux (i + 1)) 
    in aux 0 \end{whycode} \vspace*{-1.15em} 
 \begin{whygray}
  in iter (fun i -> !s = $\sum_{~0\leq j < i} a[j]$) (fun x -> s := !s + x) a; 
  !s
\end{whygray}
\end{footnotesize}
\end{frame}
\addtocounter{framenumber}{-1}


\subsubsection*{instantiate invariant}
\begin{frame}[fragile]
\frametitle{\OliveGreen{instantiate inv}}
\vspace*{0em}
\begin{footnotesize}
\begin{whycode}
let sum_iter (a: int array) 
ensures { $\texttt{result = }\sum_{~0 \leq j < \text{a.length}} a[j]$ }
= let s = ref 0 in 
  let iter ${\red \texttt{(inv:int -> prop)}}$ (f:int -> unit) (b: int array)
  requires { inv 0 } 
  ensures  { inv b.length } 
  = let rec aux (i: int)
    requires { inv i$\bwedge 0 \leq \texttt{i} \leq \texttt{b.length}$ }
    ensures  { inv b.length }
    = if i < Array.length b
      then (f b[i] ; aux (i + 1)) 
    in aux 0
  in iter ${\red \texttt{(fun i -> !s = } \sum_{~0\leq j < i} a[j])}$ (fun x -> s := !s + x) a; 
  !s
\end{whycode}
\end{footnotesize}
\end{frame}
\addtocounter{framenumber}{-1}


\begin{frame}[fragile]
\frametitle{\OliveGreen{instantiate inv}}
\vspace*{0em}
\begin{footnotesize}
\begin{whycode}
let sum_iter (a: int array) 
ensures { $\texttt{result = }\sum_{~0 \leq j < \text{a.length}} a[j]$ }
= let s = ref 0 in 
  let iter (f:int -> unit) (b: int array)
  requires { ${\red \texttt{!s = } \sum_{~0\leq j < 0} \texttt{ a[j]})}$ } 
  ensures  { ${\red \texttt{!s = } \sum_{~0\leq j < \texttt{b.length}} \texttt{a[j]}})$ } 
  = let rec aux (i: int)
    requires { inv i$\bwedge 0 \leq \texttt{i} \leq \texttt{b.length}$ }
    ensures  { inv b.length }
    = if i < Array.length b
      then (f b[i] ; aux (i + 1)) 
    in aux 0
  in iter (fun x -> s := !s + x) a; 
  !s
\end{whycode}
\end{footnotesize}
\end{frame}
\addtocounter{framenumber}{-1}



\begin{frame}[fragile]
\frametitle{\OliveGreen{instantiate inv}}
\vspace*{0em}
\begin{footnotesize}
\begin{whycode}
let sum_iter (a: int array) 
ensures { $\texttt{result = }\sum_{~0 \leq j < \text{a.length}} a[j]$ }
= let s = ref 0 in 
  let iter (f:int -> unit) (b: int array)
  requires { ${\red \texttt{!s = } \sum_{~0\leq j < 0} \texttt{ a[j]})}$ } 
  ensures  { ${\red \texttt{!s = } \sum_{~0\leq j < \texttt{b.length}} \texttt{a[j]}})$ } 
  = let rec aux (i: int)
    requires { ${\red \texttt{!s = } \sum_{~0\leq j < \texttt{i}} \texttt{a[j]}} \bwedge 0 \leq \texttt{i} \leq \texttt{b.length}$ }
    ensures  { ${\red \texttt{!s = } \sum_{~0\leq j < \texttt{b.length}} \texttt{a[j]}}$ }
    = if i < Array.length b
      then (f b[i] ; aux (i + 1)) 
    in aux 0
  in iter (fun x -> s := !s + x) a; 
  !s
\end{whycode}
\end{footnotesize}
\end{frame}
\addtocounter{framenumber}{-1}

\subsubsection*{Instantiate function f}
\begin{frame}[fragile]
\frametitle{\OliveGreen{instantiate f}}
\vspace*{0em}
\begin{footnotesize}
\begin{whycode}
let sum_iter (a: int array) 
ensures { $\texttt{result = }\sum_{~0 \leq j < \text{a.length}} a[j]$ }
= let s = ref 0 in 
  let iter ${\red \texttt{(f:int -> unit)}}$ (b: int array)
  requires { ${ \texttt{!s = } \sum_{~0\leq j < 0} \texttt{ a[j]})}$ } 
  ensures  { ${ \texttt{!s = } \sum_{~0\leq j < \texttt{b.length}} \texttt{a[j]}})$ } 
  = let rec aux (i: int)
    requires { ${ \texttt{!s = } \sum_{~0\leq j < \texttt{i}} \texttt{a[j]}} \bwedge 0 \leq \texttt{i} \leq \texttt{b.length}$ }
    ensures  { ${ \texttt{!s = } \sum_{~0\leq j < \texttt{b.length}} \texttt{a[j]}}$ }
    = if i < Array.length b
      then (${\red\texttt{f}}$ b[i] ; aux (i + 1)) 
    in aux 0
  in iter ${\red\texttt{(fun x -> s := !s + x)}}$ a; 
  !s
\end{whycode}
\end{footnotesize}
\end{frame}
\addtocounter{framenumber}{-1}

\begin{frame}[fragile]
\frametitle{\OliveGreen{instantiate f}}
\vspace*{0em}
\begin{footnotesize}
\begin{whycode}
let sum_iter (a: int array) 
ensures { $\texttt{result = }\sum_{~0 \leq j < \text{a.length}} a[j]$ }
= let s = ref 0 in 
  let iter (b: int array)
  requires { ${ \texttt{!s = } \sum_{~0\leq j < 0} \texttt{ a[j]})}$ } 
  ensures  { ${ \texttt{!s = } \sum_{~0\leq j < \texttt{b.length}} \texttt{a[j]}})$ } 
  = let rec aux (i: int)
    requires { ${ \texttt{!s = } \sum_{~0\leq j < \texttt{i}} \texttt{a[j]}} \bwedge 0 \leq \texttt{i} \leq \texttt{b.length}$ }
    ensures  { ${ \texttt{!s = } \sum_{~0\leq j < \texttt{b.length}} \texttt{a[j]}}$ }
    = if i < Array.length b
      then (${\red\texttt{s := !s + b[i]}}$ ; aux (i + 1)) 
    in aux 0
  in iter a; 
  !s
\end{whycode}
\end{footnotesize}
\end{frame}
\addtocounter{framenumber}{-1}


%\section*{Inlining}
%\begin{frame}
% \begin{itemize}
% 	\item[] $${\large \red \texttt{Inlining}}$$
% 	\begin{center}
%% 	\includegraphics[scale=0.31]{GhostBusters}
% 	\end{center} 
% \end{itemize}
%\end{frame}
%\addtocounter{framenumber}{-1}

\subsection*{our contribution}
\begin{frame}[fragile]
\frametitle{our contribution}
\begin{displaymath}
\begin{scriptsize}
\begin{array}{@{\hspace*{-2.5em}}lllllllllll}
  \dfrac{\texttt{HO program}}{\texttt{ + spec}} & \texttt{->} 
& \fbox{\red Inlining}& \texttt{->} & \dfrac{\texttt{FO program}}{\texttt{ + spec}} & \texttt{->} 
& \fbox{WP Calculus} & \texttt{->} & \texttt{VCs} & \texttt{->} 
& \fbox{ATPs}
\end{array}
\end{scriptsize}
\end{displaymath}\pause
\begin{enumerate}
\item identify the fragment of ML where inlining can be done
\begin{itemize}
\item - syntax, semantics, type system with effects, typing restrictions
\item - formal description of inlining  \\[1em]
\end{itemize} \pause
\item proof of inlining correctness
\begin{itemize}	
\item termination proof
\item semantics preservation proof
\end{itemize}
\end{enumerate}
\end{frame}



% 2. quelles difficultés
%    ce qui nous amène à Mini-ML = ML + références globales
%    syntaxe, sémantique, système de types avec effets, preuve de type soundness
%    identification d'un fragment de Mini-ML + définition inlining + preuve

%\begin{frame}\frametitle<1->{\quad Generalizing approach for second-order programs}
%%\frametitle<2->{\quad \color{red} }
%\vspace*{-4em}
%An annotated {\red second-order} program $\mathcal{P}$ of some language $\mathcal{L}$ with 
%\begin{itemize}
%\item[-]<1-> data structures: arrays, second-order functions, ...
%\item[-]<1-> side effects: mutable references, recursive functions, ...
%\end{itemize}
%becomes a {\red first-order} annotated program $\mathcal{I(P)}$.\\[2em]
%
%\only<2->{\textbf{Can we use the VCs generated from $\mathcal{I(P)}$ to verify  $\mathcal{P}$}? } ~ \\[0.5em]
%\only<4->{{\blue \quad Yes, but we must prove that $\mathcal{P}$ and $\mathcal{I(P)}$ are equivalent.}} ~ \\[0.5em]
%\only<3->{\textbf{ Can we generalize our approach for \textit{any} such program $\mathcal{P}$?}} ~ \\[0.5em]
%\only<5->{{\blue \quad It depends on what is $\mathcal{L}$.}} 
%%\only<5->{{\blue \quad No for any \textit{arbitrary} second-order program.} } ~ \\
%%\only<5->{{\blue \quad Yes for any program of some restricted \textit{subset} of programs similar to \texttt{sum\_iter}}}.
%\end{frame}

\subsection*{identify fragment}
\begin{frame}[fragile]\frametitle{les liaisons dangereuses}
\begin{small}
\begin{whycode} 
${\red \only<1>{\mathcal{P}_1}\only<2->{\mathcal{P}_2}}$: 
   let F = 
      ($\only<1>{\texttt{\red r:= !r + 1;}}\only<2->{\texttt{\red diverge(); }}$ fun (f:int -> int) -> f !r)
   in $\only<1>{\texttt{(F id) + (F id)}}\only<2->{\texttt{()}}$
   
*?%$\hspace*{8em}\only<2>{{\red \mathcal{P}_{1_{\{r=0\}}} \longrightarrow^{\star} 2_{\{ r=1\}}}}\only<4>{\red \mathcal{P}_{2_{\{r=0\}}} \longrightarrow^{\star} \infty}$
?* 
${\red\only<1>{ \mathcal{I}(P_1)}\only<2->{\mathcal{I}(P_2)}}: \only<1>{\texttt{(r:= !r + 1; !r) + (r:= !r + 1; !r)}}
\only<2->{\texttt{()}}$
   
*?%$\hspace*{8em}\only<2>{\mathcal{I}(P_1)_{\{r=0\}} \longrightarrow^{\star} 3_{\{r=2\}}}\only<4>{\mathcal{I}(P_2)_{\{r=0\}} \longrightarrow^{\star} ()_{\{r=0\}}}$      
?*
\end{whycode}
\end{small}
\end{frame}

\subsection*{type system with effects}
\begin{frame}[fragile]
\frametitle{experimental language}
%\hspace*{-2em} \mbox{{\blue Our approach}: identify the subset of ML language in which we can write}\vspace*{-0.5em}
%\begin{itemize}
%\item[-] interesting examples like sum\_iter;
%\item[-] only programs we can inline correctly.\\[1em]
%\end{itemize}\pause
%{\hspace*{-2em} \blue A hint of the solution}: 
a fragment of ML with
\begin{itemize}
	\item[-] references, recursion, integers, integer lists
	\item[-] {\red second-order} programs with types \\[0.5em]
		\begin{itemize}	
		\begin{footnotesize} 
			\item[] \qquad$\tau^{\red 0}:\{\texttt{ int, bool,... }\}$
			\item[] \qquad$\tau^{\red 1}:\{\texttt{ int -> bool, int -> (int -> bool), ... }\}$
			\item[] \qquad$\tau^{\red 2}:\{\texttt{ (int -> int) -> bool, ... }\}$\\[1em]
		\end{footnotesize}
		\end{itemize}
		 \pause
	\item[-] a type system {\red with effects}:\\[0.5em]
		\begin{itemize}
			\item[] \qquad $\Gamma \vdash t: \tau^{i}, {\red \theta, \rho} $
		\end{itemize}
\end{itemize}
\end{frame}


\subsection*{Example: typing of local bindings}
\begin{frame}[fragile]
\frametitle{les liaisons moins dangereuses}
\begin{small}
%\begin{adjustwidth}{-1.5em}{0em}
%\begin{displaymath}
%		\dfrac
%		 {\Gamma_1 \vdash t_1 : \tau^i, \theta_1, \rho_1
%		  \qquad \Gamma_2, x:\tau^i  \vdash t_2 : \tau^j, \theta_2, \rho_2 
%		  \qquad \only<1>{\phantom{\textcolor{red}{(i = 2) \Rightarrow 
%			\boldsymbol{(\theta_1 = \bot_{\theta}\bwedge \rho_1 = \bot_{\rho})}}}}			
%			\only<2->{ \textcolor{red}{(i = 2) \Rightarrow 
%			\boldsymbol{(\theta_1 = \bot_{\theta}\bwedge \rho_1 = \bot_{\rho})}}}}
%		 {\Gamma_1 \vdash \text{ let } x = t_1 \text{ in } t_2 
%				: \tau^{j}, 
%				(\theta_1 \bvee \theta_2), 
%				(\rho_1 \bvee \rho_2)}
%\end{displaymath}		
%\end{adjustwidth}	

		\infrule[] 
			{\vspace*{0.5em}\hspace*{-2.1em}\Gamma \vdash t_1 : \tau^i, \Rouge{\theta_1}, \Rouge{\rho_1} 
			\quad \Gamma, x:\tau^i  \vdash t_2 : \tau^j, \Rouge{\theta_2}, \Rouge{\rho_2} 
			\quad 
			\only<1>{\phantom{\textcolor{blue}{(i = 2) \Rightarrow 
			\boldsymbol{(\theta_1 = \bot_{\theta}\bwedge \rho_1 = \bot_{\rho})}}}}			
			\only<2->{ \textcolor{blue}{(i = 2) \Rightarrow 
			(\theta_1 = \bot_{\theta}\bwedge \rho_1 = \bot_{\rho})}}} 
			{\vspace*{-1em}\hspace*{-2.1em}\Gamma \vdash \text{ let } x = t_1 \text{ in } t_2 
				: \tau^{j}, 
				(\Rouge{\theta_1 \bvee \theta_2}), 
				(\Rouge{\rho_1 \bvee \rho_2})}
	\end{small}
	%\end{adjustwidth}	
\end{frame}

\section*{Ghost Code}

\begin{frame}
 \begin{itemize}
 	\item[] $${\large \red \texttt{Ghost Code}}$$
 	\begin{center}
 	\includegraphics[scale=0.31]{GhostBusters}
 	\end{center} 
 \end{itemize}
\end{frame}
\addtocounter{framenumber}{-1}

\subsection*{ghost properties informally}
% un pgm avec du code ghost :
%  - pour le prouver, je prends *tout* le code, sans distinction
%  - pour l'exécuter, j'oublie le code ghost

\begin{frame}
\frametitle{\red{ghost code}}
	may be inserted for the purpose of specification and/or proof. \\[1em]
	\pause	
	\begin{itemize}
	\item[\red $\bullet$] regular code does not see ghost data
	\item[\red $\bullet$] ghost code may read regular data (but cannot modify it)\\[1em]
	\end{itemize}
in particular, ghost code may be removed without observable modification
\end{frame}


\subsection*{example: sum of list elements} 
\begin{frame}[fragile]
\frametitle{the sum of list elements}
\begin{footnotesize}
\begin{whycode}
 let sum_list (l: int list) : int
 ensures { result = $\sum$ l }
 = let s = ref 0 in
   List.iter 
     (${\red\texttt{fun before -> !s = }\sum \texttt{ before}}$) (fun x -> s := !s + x) l; 
   !s     
\end{whycode}
\end{footnotesize}
\end{frame}

\begin{frame}[fragile]
\frametitle{precondition and postcondition}
\begin{footnotesize}
\begin{whycode}
let iter ($\textcolor{red}{\text{inv:int list -> prop}}$) (f:int -> unit) (l: int list)
requires { ${\red \texttt{inv Nil}}$ }
ensures  { ${\red \texttt{inv (reverse l)}}$ } 
= ${\gray\texttt{let rec loop (before:int list) (after:int list)}}$
  ${\gray\texttt{requires \{ inv (before) \} }}$
  ${\gray\texttt{ensures~ \{ inv (rev\_append after before) \}}}$  
  ${\gray\texttt{= match after with}}$
     ${\gray\texttt{| Nil~~~~~~ -> ()}}$
     ${\gray\texttt{| Cons x xs -> f x; loop (Cons x before) xs}}$
  ${\gray\texttt{in loop Nil l }}$
\end{whycode}
\end{footnotesize}
\end{frame}
\addtocounter{framenumber}{-1}



\begin{frame}[fragile]
\frametitle{loop invariant}
\begin{footnotesize}
\begin{whycode}
${\gray\texttt{let iter (inv:int list -> prop)}} ~
{\gray\texttt{(f:int -> unit) (l:int list)}}$
${\gray\texttt{requires \{ inv Nil \}}}$
${\gray\texttt{ensures~ \{ inv (reverse l)\}}}$    
= let rec loop $(\textcolor{red}{\text{before:int list}})$ (after: int list) 
  requires $\texttt{\{ inv {\red (before)} \}}$
  ensures  $\texttt{\{ inv {\red(rev\_append after before)} }\}$  
  ${\gray\texttt{= match after with}}$
     ${\gray\texttt{| Nil~~~~~~ -> ()}}$
     ${\gray\texttt{| Cons x xs -> f x; loop (Cons x before) xs}}$
  ${\gray\texttt{in loop Nil l }}$
\end{whycode}
\end{footnotesize}
\end{frame}
\addtocounter{framenumber}{-1}


\begin{frame}[fragile]
\frametitle{regular code}
\begin{footnotesize}
\begin{whycode}
${\gray\texttt{let iter (inv:int list -> prop)}} ~
{\gray\texttt{(f:int -> unit) (l:int list)}}$
${\gray\texttt{requires \{ inv Nil \}}}$
${\gray\texttt{ensures~ \{ inv (reverse l)\}}}$    
= ${\gray\texttt{let rec loop (before:int list) (after:int list)}}$
  ${\gray\texttt{requires \{ inv (before) \} }}$
  ${\gray\texttt{ensures~ \{ inv (rev\_append after before) \}}}$
  = match after with 
     | Nil       -> ()
     | Cons x xs -> f x; loop ($\textcolor{red}{\text{Cons x before}}$) xs 
  in loop $\texttt{\red Nil}$ l 
\end{whycode}
\end{footnotesize}
\end{frame}
\addtocounter{framenumber}{-1}

\begin{frame}[fragile]
\frametitle{\red{ghost code}}
\begin{footnotesize}
\begin{whycode}
let iter ($\textcolor{red}{\text{ghost inv:int list -> prop}}$) (f:int -> unit) (l: int list)
requires { ${ \texttt{inv Nil}}$ }
ensures  { ${ \texttt{inv (reverse l)}}$ }
= let rec loop $(\textcolor{red}{\text{ghost before:int list}})$ (after: int list) 
  requires $\texttt{\{ inv { (before)} \}}$
  ensures  $\texttt{\{ inv {(rev\_append after before)} }\}$     
  = match after with 
     | Nil       -> ()
     | Cons x xs -> f x; loop ($\textcolor{red}{\text{Cons x before}}$) xs 
  in loop $\texttt{\red Nil}$ l 
\end{whycode}
\end{footnotesize}
\end{frame}
\addtocounter{framenumber}{-1}


\subsection*{Ghost formalisation}

% 4. formalisation de Ghost-ML
%    syntaxe, sémantique, typage avec effets (montrer le jugement),
%      bien expliquer qu'il y a des choix à faire (mais aussi des*)
%      obligations), l'opération erasure et la preuve

\begin{frame}
\frametitle{our contribution}
\begin{itemize}
\item formalize Ghost-ML: syntax, semantics, type system with ghost effects \pause
\item prove ghost code non-interference with regular code \pause
\item define ghost erasure and prove semantics preservation  
\end{itemize}
\end{frame}

\begin{frame}
\frametitle{ghost code non-interference}
type system with effects 
$$\Gamma \vdash_{gh} t: \tau^{i}, \theta, \rho, {\red \beta} $$

\bigskip\pause
typing invariant 
$$\Gamma \vdash_{gh} t : \tau, (\theta, \rho, \gbr) 
 ~~ \Rightarrow ~~
\models (\gbr = \gbtr) \Rightarrow (\theta = \bth \bwedge \rho = \brh) $$
\end{frame}


\begin{frame}
\frametitle{ghost code propagation}
\begin{footnotesize}	
\begin{displaymath}
\begin{array}{l}
\dfrac
	{
		\ghosttyping
			{t_1}
			{\tau_1}
			{\theta_1}
			{\rho_1}
			{\gbra{1}^{\rouge{'}}}
			{} \qquad
		\ghosttyping
			{t_2}
			{\tau_2}
			{\theta_2}
			{\rho_2}
			{\gbra{2}}
			{  \quad \models \gba{1} \Rightarrow \gba{1}^{'}}
	}
	{\ghosttyping
		{let~x^{\gbra{1}}_{\tau_1} = t_1~in~t_2}
		{\tau_2}
		{\theta_1 \bvee \theta_2}
		{\rho_1 \bvee \rho_2}
		{\gbra{2} \bvee (\gbran{1} \bwedge \gbra{1}^{\rouge{'}})}
		{}}
	{}
	\end{array}	
\end{displaymath}
\end{footnotesize}	
\end{frame}

\section*{Conclusion}
\begin{frame}
\frametitle{conclusion}
\begin{itemize}
\item experimental evaluation 
\begin{itemize}
	\item inlining can be useful 
\end{itemize}
\item formalization
\begin{itemize}
	\item inlining, ghost-code, correctness proofs
\end{itemize}
\end{itemize}
\end{frame}
%\begin{frame}\frametitle{To summarize, in our work}\vspace*{-4em}
%\begin{itemize}
%	\item[\blue$-$]  we described of the inlining by the set of rewriting rules; \\[1em]
%	\item[\blue$-$]  we proved that the procedure described by these rules \\[0.5em]
%	
%	\begin{itemize}
%		\item[$\blacktriangleright$]
%		{\normalsize always {\red terminates } giving a first-order output term} \\[0.3em]
%		\item[$\blacktriangleright$]  
%		{\normalsize {\red preserves} the semantics of the initial term} \\[1em]
%	\end{itemize}
%	\item[\blue $-$] we formalized the notion of {\red ghost code}
%\end{itemize}
%\end{frame}

	
% In this work (1) formalize domain 
% prove soundness 

% Formaliziation : typing
% Correctness : equivalence
% Example
\end{document}

% Local Variables:
% compile-command: "rubber -d slides"
% ispell-local-dictionary: "francais"
% End:



